\documentclass[12pt,a4paper]{article}
\usepackage[utf8]{inputenc}
\usepackage[T1]{fontenc}
\usepackage[polish]{babel}
\usepackage{geometry}
\usepackage{enumitem}
\usepackage{titlesec}
\usepackage{parskip}
\usepackage{fancyhdr}
\usepackage{xcolor}
\usepackage{mdframed}
\usepackage{microtype}
\usepackage{booktabs}
\usepackage{hyperref}
\usepackage{multirow}
\usepackage{array}
\usepackage{longtable}
\usepackage{xfp}

\geometry{margin=2.5cm}

\hypersetup{colorlinks=true, linkcolor=black, urlcolor=blue}

\pagestyle{fancy}
\fancyhf{}
\rhead{SDLC-LLM Benchmark}
\lhead{Zadanie R3 -- Conflict Detection}
\rfoot{\thepage}

\titleformat{\section}{\large\bfseries}{}{0em}{}[\titlerule]
\titleformat{\subsection}{\normalsize\bfseries}{}{0em}{}
\titleformat{\subsubsection}{\normalsize\itshape}{}{0em}{}

\definecolor{metabg}{gray}{0.93}
\definecolor{promptbg}{RGB}{235,245,255}
\definecolor{gptcolor}{RGB}{16,163,127}
\definecolor{claudecolor}{RGB}{210,105,30}
\definecolor{llamacolor}{RGB}{100,100,180}
\definecolor{geminicolor}{RGB}{0,200,180}
\definecolor{score3}{RGB}{220,240,220}
\definecolor{score2}{RGB}{255,243,205}
\definecolor{score1}{RGB}{255,220,200}
\definecolor{score0}{RGB}{255,200,200}

\newcommand{\scorebox}[1]{%
  \ifnum#1=3\colorbox{score3}{\textbf{3}}%
  \else\ifnum#1=2\colorbox{score2}{\textbf{2}}%
  \else\ifnum#1=1\colorbox{score1}{\textbf{1}}%
  \else\colorbox{score0}{\textbf{0}}\fi\fi\fi
}

\newcommand{\fonescorebox}[1]{%
    \ifnum\fpeval{#1 >= 0.8} = 1
        \colorbox{score3}{\textbf{#1}}% <--- TUTAJ JEST ZMIANA
    \else
        \ifnum\fpeval{#1 >= 0.6} = 1
            \colorbox{score2}{\textbf{#1}}% <--- I TUTAJ
        \else
            \ifnum\fpeval{#1 >= 0.4} = 1
                \colorbox{score1}{\textbf{#1}}% <--- I TUTAJ
            \else
                \colorbox{score0}{\textbf{#1}}% <--- I TUTAJ
            \fi
        \fi
    \fi
}

\begin{document}

% ─── STRONA TYTUŁOWA ────────────────────────────────────────────────
\begin{center}
  {\LARGE\bfseries Zadanie R3 -- Conflict Detection}\\[0.5em]
  {\large Wykrywanie i rozwiązywanie konfliktów logicznych w opisie systemu}\\[1em]
  {\large Case study: Appointment Booking System}\\[0.3em]
  {\large Tryb: LLM-only \quad }\\[0.3em]
  {\large Modele: \texttt{gpt-5.4} \quad \texttt{claude-sonnet-4-6} \quad \texttt{llama-3.1-8b}} \quad \texttt{gemini-3-flash-preview}}
\end{center}

\vspace{1em}\hrule\vspace{1.5em}

% ─── PROMPT ─────────────────────────────────────────────────────────
\section{Prompt}

\begin{mdframed}[backgroundcolor=metabg, linewidth=0pt, innerleftmargin=10pt, innerrightmargin=10pt, innertopmargin=8pt, innerbottommargin=8pt]
\textbf{System prompt:}\\
\small Jesteś Doświadczonym Analitykiem Wymagań (Senior Business Analyst). Twoją specjalnością jest krytyczna weryfikacja surowych koncepcji biznesowych. Twoim zadaniem jest dogłębna analiza dostarczonego opisu systemu pod kątem wykrywania konfliktów (sytuacji, w których reguła A wyklucza regułę B lub proces nie może zostać ukończony z powodu sprzecznych założeń). Rozwiązując konflikty, operuj wyłącznie w ramach dostarczonego kontekstu. Nie dodawaj nowych, niewymienionych wcześniej modułów ani funkcjonalności.
\end{mdframed}

\vspace{0.5em}
\begin{mdframed}[backgroundcolor=promptbg, linewidth=0pt, innerleftmargin=10pt, innerrightmargin=10pt, innertopmargin=8pt, innerbottommargin=8pt]
\textbf{User prompt (struktura):}\\
\small
Przeanalizuj poniższy surowy opis systemu w tagach \texttt{<opis>}. Zwróć wynik w formacie Markdown z dwiema sekcjami: \textbf{(1) Lista wykrytych konfliktów} -- każdy punkt zaczyna się od \texttt{- }, bez numeracji; \textbf{(2) Poprawiona wersja} -- tekst ciągły bez wypunktowań.
\end{mdframed}

\vspace{0.5em}
\begin{mdframed}[backgroundcolor=metabg, linewidth=0pt, innerleftmargin=10pt, innerrightmargin=10pt, innertopmargin=8pt, innerbottommargin=8pt]
\textbf{Wejście -- opis systemu (celowo sprzeczny):}\\[0.3em]
\small
System rezerwacji wizyt ma na celu ułatwienie pacjentom (rola: User) umawiania spotkań ze specjalistami (rola: Specialist), podczas gdy nad całością infrastruktury czuwa administrator (rola: Admin). Głównym założeniem biznesowym jest to, że każdy użytkownik może rezerwować dowolną, nielimitowaną liczbę terminów, pod warunkiem jednak, że w danym momencie nie posiada więcej niż 3 aktywnych rezerwacji. Do użytkownika może zostać przypisana jedna z dwóch ról: Secjalista, User.
Ze względu na bezpieczeństwo dostępu do informacji, logi systemu nie są zapisywane. Operacja rezerwacji musi zostać zrealizowana w czasie mniejszym niż jedna sekunda w momencie standardowego obłożenia systemu.

Platforma gwarantuje całkowity brak podwójnych rezerwacji (tzw. double booking) – system jest zaprojektowany tak, że jeden, sztywny, 30-minutowy slot czasowy może zostać zajęty wyłącznie przez jedną osobę (tylko jeżeli termin jest dostępny). Jednocześnie, aby zmaksymalizować obłożenie placówki, w przypadku najbardziej obleganych lekarzy system dopuszcza tzw. overbooking, pozwalając kilku różnym pacjentom zarezerwować dokładnie ten sam termin. Terminy specjalistów, które są oznaczone jako zajęte mogą zostać zarezerwowane jeszcze przez jednego użytkownika tylko w momecie, w którym dany specjalista nie już dostępnych terminów.

Zarządzanie harmonogramem leży całkowicie w gestii specjalistów. Mają oni absolutną swobodę zarządzania i mogą usunąć ze swojego kalendarza dowolny, widoczny w systemie termin, wliczając w to wizyty, które zostały już potwierdzone i zarezerwowane przez pacjentów. System dba o spójność danych i bezwzględnie zabrania specjalistom usuwania jakichkolwiek okienek czasowych, jeśli przypisano już do nich użytkownika.

Kwestia rezygnacji z wizyty została zaprojektowana z myślą o maksymalnej wygodzie: użytkownik ma prawo anulować swoją rezerwację w dowolnym momencie, wtedy status rezerwacji zmienia się na CANCELLED, a termin specjalisty w którym ta rezerwacja miała mieć miejsce zostaje niedostępny. Wizyty oznaczone statusem CANCELLED zwalniają termin dostępności u specjalistów. Należy pamiętać, że zgodnie z regulaminem, aplikacja zablokuje możliwość anulowania, jeśli do spotkania pozostało mniej niż 24 godziny. Po pomyślnym anulowaniu wizyty przez pacjenta, dany slot czasowy zostaje bezpowrotnie usunięty z kalendarza. Wraca on do ogólnodostępnej puli, aby inni pacjenci mogli z niego skorzystać. Z kolei sama wizyta, po przejściu przez system anulacji, otrzymuje jeden z dostępnych w cyklu życia statusów: AVAILABLE, BOOKED lub COMPLETED. Każda rezerwacja oraz anulowanie generuje logi, które są trwale zapisywane w systemie. System natomiast nie posiada sprecyzowanych wymagań dotyczących czasu odpowiedzi dla użytkownika ze względu na niewielkie zasoby obliczeniowe.

\end{mdframed}

% ─── WYNIKI PER MODEL ───────────────────────────────────────────────
\section{Wyniki}

% ════════════════════════════════════════════
\subsection{\textcolor{gptcolor}{GPT-5.4} \quad {\normalfont\small czas odpowiedzi: 30.7s}}

\subsubsection{Lista wykrytych konfliktów}
\begin{itemize}[leftmargin=1.5em, itemsep=2pt]
  \item Jednoczesna gwarancja braku double booking i dopuszczenie overbookingu dla tego samego terminu.
  \item Absolutna swoboda specjalisty usuwania terminów vs.\ bezwzględny zakaz usuwania slotów z przypisanym użytkownikiem.
  \item Możliwość anulowania ``w dowolnym momencie'' vs.\ blokada anulowania na mniej niż 24h przed wizytą.
  \item Slot po anulowaniu ``bezpowrotnie usunięty'' i jednocześnie ``wraca do puli'' -- wzajemnie wykluczające się skutki.
  \item Status CANCELLED opisany jako wynik anulowania, ale nieujęty w cyklu życia (AVAILABLE, BOOKED, COMPLETED).
  \item Możliwość usuwania przez specjalistę zarezerwowanych wizyt stoi w sprzeczności z deklaracją spójności danych.
\end{itemize}

\subsubsection{Poprawiona wersja (podsumowanie)}
\small
Każdy użytkownik może rezerwować kolejne terminy, jednak nie może posiadać więcej niż 3 aktywnych rezerwacji. Platforma działa w oparciu o sztywne 30-minutowe sloty i nie dopuszcza podwójnych rezerwacji. Specjalista może usuwać wyłącznie wolne terminy -- slotów z istniejącą rezerwacją usunąć nie można. Użytkownik może anulować rezerwację wyłącznie do momentu, gdy do spotkania pozostało co najmniej 24h; po skutecznym anulowaniu status zmienia się na CANCELLED, a slot wraca do puli. Cykl życia: AVAILABLE, BOOKED, COMPLETED, CANCELLED.
\normalsize

\vspace{0.5em}
\begin{tabular}{lccl}
\toprule
\textbf{Metryka} & \textbf{Ocena} & & \textbf{Uzasadnienie} \\
\midrule
M1 Correctness    & \scorebox{3} && Wszystkie 6 konfliktów wykryte poprawnie, poprawiona wersja spójna. \\
M2 Completeness   & \scorebox{3} && Żaden z kluczowych konfliktów nie został pominięty. \\
M3 Consistency    & \scorebox{3} && Poprawiona wersja eliminuje wszystkie wykryte sprzeczności. \\
M4 Clarity        & \scorebox{3} && Konflikty opisane zwięźle, poprawiona wersja czytelna. \\
M5 Maintainability & \scorebox{3} && Wynik gotowy do użycia jako wejście dla kolejnego etapu. \\
\bottomrule
\end{tabular}

\vspace{1.5em}

% ════════════════════════════════════════════
\subsection{\textcolor{claudecolor}{Claude Sonnet 4.6} \quad {\normalfont\small czas odpowiedzi: 16.54s}}

\subsubsection{Lista wykrytych konfliktów}
\begin{itemize}[leftmargin=1.5em, itemsep=2pt]
  \item Gwarancja braku double booking jednocześnie z dopuszczeniem overbookingu -- bezpośrednia sprzeczność logiczna.
  \item Absolutna swoboda specjalisty vs.\ bezwzględny zakaz usuwania slotów z przypisanym użytkownikiem.
  \item Slot po anulowaniu ``bezpowrotnie usunięty'' i jednocześnie ``wraca do puli''.
  \item Status CANCELLED wynika z anulowania, ale nie figuruje wśród statusów cyklu życia (AVAILABLE, BOOKED, COMPLETED).
  \item ``Nielimitowana liczba terminów'' sprzeczna z twardym limitem 3 aktywnych rezerwacji.
\end{itemize}

\subsubsection{Poprawiona wersja (podsumowanie)}
\small
Każdy użytkownik może rezerwować kolejne terminy bez odgórnego limitu łącznej liczby wizyt, pod warunkiem że w danym momencie nie posiada więcej niż 3 aktywnych rezerwacji o statusie BOOKED. System nie dopuszcza overbookingu -- jeden 30-minutowy slot może zajmować wyłącznie jeden pacjent. Specjalista może usunąć lub zmodyfikować wyłącznie sloty bez przypisanego pacjenta. Anulowanie jest możliwe tylko gdy do wizyty pozostało co najmniej 24h; po anulowaniu slot wraca do puli ze statusem AVAILABLE. Cykl życia: AVAILABLE, BOOKED, CANCELLED, COMPLETED.
\normalsize

\vspace{0.5em}
\begin{tabular}{lccl}
\toprule
\textbf{Metryka} & \textbf{Ocena} & & \textbf{Uzasadnienie} \\
\midrule
M1 Correctness    & \scorebox{3} && Wykryte konflikty merytorycznie poprawne, w tym subtelny konflikt ``nielimitowana vs.\ limit 3''. \\
M2 Completeness   & \scorebox{2} && Pominięty konflikt dot.\ usuwania zarezerwowanych wizyt przez specjalistę (6. punkt GPT). \\
M3 Consistency    & \scorebox{3} && Poprawiona wersja w pełni spójna, dobrze wyjaśnia statusy cyklu życia. \\
M4 Clarity        & \scorebox{3} && Najczytelniejsza z trzech odpowiedzi, cykl życia opisany z definicjami. \\
M5 Maintainability & \scorebox{3} && Wynik gotowy do użycia, poprawiona wersja szczegółowa i precyzyjna. \\
\bottomrule
\end{tabular}

\vspace{1.5em}

% ════════════════════════════════════════════
\subsection{\textcolor{llamacolor}{Llama 3.1 8B} \quad {\normalfont\small czas odpowiedzi: 1.88s}}

\subsubsection{Lista wykrytych konfliktów}
\begin{itemize}[leftmargin=1.5em, itemsep=2pt]
  \item Co dzieje się ze slotem gdy użytkownik anuluje rezerwację -- brak wyjaśnienia. \textit{(niejednoznaczność, nie konflikt)}
  \item Co dzieje się ze slotem gdy użytkownik nie przychodzi i nie anuluje. \textit{(niejednoznaczność, nie konflikt)}
  \item Czy anulowanie w ciągu 24h powinno być możliwe. \textit{(pytanie, nie wykryty konflikt)}
  \item Co dzieje się gdy użytkownik próbuje przekroczyć limit 3 rezerwacji. \textit{(niejednoznaczność, nie konflikt)}
  \item Co dzieje się gdy specjalista usuwa kilka rezerwacji naraz. \textit{(edge case, nie konflikt)}
\end{itemize}

\textit{Uwaga: model nie wykrył żadnego rzeczywistego konfliktu logicznego. Zamiast tego sformułował pytania o niejednoznaczności i edge cases. Kluczowe sprzeczności (double booking vs.\ overbooking, slot usunięty vs.\ wraca do puli, CANCELLED poza cyklem życia) zostały całkowicie pominięte.}

\subsubsection{Poprawiona wersja (podsumowanie)}
\small
Model w sekcji 2 przepisał opis niemal identycznie z wejściem, zachowując wszystkie oryginalne konflikty (w tym overbooking, absolutną swobodę specjalisty, slot bezpowrotnie usunięty). Dodał jedynie zdanie o możliwości usunięcia slotu przez specjalistę po anulowaniu przez użytkownika lub niestawieniu się pacjenta -- co stanowi nową funkcjonalność, której dodawania prompt zabraniał.
\normalsize

\vspace{0.5em}
\begin{tabular}{lccl}
\toprule
\textbf{Metryka} & \textbf{Ocena} & & \textbf{Uzasadnienie} \\
\midrule
M1 Correctness    & \scorebox{0} && Model nie wykrył żadnego rzeczywistego konfliktu -- wszystkie punkty to pytania o niejednoznaczności. \\
M2 Completeness   & \scorebox{0} && Żaden z 5--6 kluczowych konfliktów nie został zidentyfikowany. \\
M3 Consistency    & \scorebox{1} && Poprawiona wersja zachowuje oryginalne konflikty, dodaje nową funkcjonalność wbrew instrukcji. \\
M4 Clarity        & \scorebox{1} && Struktura formalna zachowana, ale treść nie realizuje celu zadania. \\
M5 Maintainability & \scorebox{0} && Wynik bezużyteczny jako wejście dla kolejnego etapu -- konflikty pozostają nierozwiązane. \\
\bottomrule
\end{tabular}


% ════════════════════════════════════════════
\subsection{\textcolor{geminicolor}{Gemini 3 flash preview} \quad {\normalfont\small czas odpowiedzi: 13.96s}}

\subsubsection{Lista wykrytych konfliktów}
\begin{itemize}
    \item[-] Poprawne wskazanie konfliktu ról, w opisie została zawarta informacja, że istnieją 2 dostępne role dla użytkowników, natomiast zostały opisane 3.

    \item[-] Model słusznie zwrócił uwagę na niespojność w kontekście zapisywania logów z rezerwacji oraz anulowania wizyt.

    \item[-] Słusznie została wypunktowana niespójność w kontekście wymagań systemowych dotyczących wydajności platformy.
    
    \item[-] Trafnie zidentyfikowana sprzeczność, która prowadziła do overbookingu dla niektórych lekarzy.
    
    \item[-] Model zwrócił uwagę, że przyznanie specjalistom absolutnej swobody w usuwaniu dowolnych terminów koliduje z bezwzględnym zakazem usuwania okienek czasowych, do których przypisano już użytkownika. 
    
    \item[-] Model wykrył powszechny konflikt wymagań, w którym obietnica pełnej elastyczności dla użytkownika jest negowana przez twardy warunek brzegowy (24 godziny).

    \item[-] Odpowiedź modelu zawiera wychwycony błąd w definicji maszyny stanów; status przypisywany w jednym procesie (CANCELLED) nie istnieje w ogólnej definicji cyklu życia wizyty.
    
    \item[-] Informacja o bezpowrotnym usunięciu slotu z kalendarza po anulowaniu wyklucza się z założeniem, że wraca on do ogólnodostępnej puli dla innych pacjentów
    
\end{itemize}


\subsubsection{Poprawiona wersja (podsumowanie)}
\small
Model sprawnie wyeliminował błędy logiczne, priorytetyzując spójność danych i standardy określone w dokumentacji wzorcowej. Rozstrzygnął konflikt ról, zachowując podział na Usera, Specialistę i Admina , oraz ujednolicił zasady rezerwacji, utrzymując limit trzech aktywnych wizyt i całkowicie odrzucając overbooking na rzecz sztywnych slotów. Sprzeczne uprawnienia specjalisty rozwiązano przez zakaz usuwania zarezerwowanych terminów , a politykę anulacji ograniczono do okna powyżej 24 godzin przed wizytą. Model uporządkował cykl życia wizyty, włączając status CANCELLED obok AVAILABLE, BOOKED i COMPLETED , decydując jednocześnie o powrocie slotu do puli po rezygnacji. Ostatecznie wybrano wariant z trwałym zapisem logów dla celów audytu oraz zachowano rygorystyczny wymóg wydajności poniżej jednej sekundy. Nie sprecyzował jednak kluczowej informacji dotyczącej ilości rezerwowanych terminów przez jednego użytkownika, co może prowadzić do sztucznych rezerwacji, a następnie anulowania wizyt. Nie zwrócił także uwagi na możliwość rezerwacji niedostępnych terminów. 
\normalsize

\vspace{0.5em}
\begin{tabular}{lccl}
\toprule
\textbf{Metryka} & \textbf{Ocena} & & \textbf{Uzasadnienie} \\
\midrule
M1 Correctness    & \scorebox{2} && Model wykrył kluczowe sprzeczności zawarte w opisie systemu. \\
R1 F1-score   & \fonescorebox{0.89} && Wartość zarówno współczynnika precision jak i recall jest równe 1. \\
M3 Consistency    & \scorebox{3} && Poprawiona wersja jest w pełni spójna i logicznie poprawna. \\
\bottomrule
\end{tabular}


% ─── TABELA PORÓWNAWCZA ─────────────────────────────────────────────
\section{Porównanie modeli}

\begin{tabular}{lccccc|c}
\toprule
\textbf{Model} & \textbf{M1} & \textbf{R1} & \textbf{M3} & \textbf{M4} & \textbf{M5} & \textbf{Suma} \\
\midrule
GPT-5.4            & \scorebox{3} & \scorebox{3} & \scorebox{3} & \scorebox{3} & \scorebox{3} & \textbf{15} \\
Claude Sonnet 4.6  & \scorebox{3} & \scorebox{2} & \scorebox{3} & \scorebox{3} & \scorebox{3} & \textbf{14} \\
Llama 3.1 8B       & \scorebox{0} & \scorebox{0} & \scorebox{1} & \scorebox{1} & \scorebox{0} & \textbf{2}  \\
Gemini 3 Flash Preview & \scorebox{2} & \fonescorebox{0.89} & \scorebox{3} & \scorebox{0} & \scorebox{0} & \textbf{5.89} \\
\bottomrule
\end{tabular}

\vspace{0.5em}
\begin{mdframed}[backgroundcolor=metabg, linewidth=0pt, innerleftmargin=10pt, innerrightmargin=10pt, innertopmargin=6pt, innerbottommargin=6pt]
\small
\textbf{Czas odpowiedzi:} GPT-5.4: 30.7s \quad Claude: 16.54s \quad Llama: 1.88s \quad Gemini: 13.96s\\
Llama odpowiedziała najszybciej, ale wynik jest bezużyteczny. Zarówno Gemini jak i GPT osiągnęły porównywalną jakość oraz wykryli maksymalną liczbę konfliktów, a Claude nie zidentyfikował jednej sprzeczności, jednakże lepiej opisał cykl życia statusów.
\end{mdframed}

\end{document}